\documentclass[aspectratio=169,10pt]{beamer}

% ---------- Theme ----------
\usepackage{beamerthemeAcademic}
\useblue

% ---------- Fonts ----------
\usepackage{fontspec}
\usepackage{xeCJK}
\IfFontExistsTF{Microsoft YaHei}{
  \setCJKmainfont{Microsoft YaHei}
  \setCJKsansfont{Microsoft YaHei}
}{
  \IfFontExistsTF{SimSun}{
    \setCJKmainfont{SimSun}
    \setCJKsansfont{Microsoft YaHei}
  }{
    \setCJKmainfont{Noto Serif CJK SC}
    \setCJKsansfont{Noto Sans CJK SC}
  }
}
\IfFontExistsTF{Consolas}{\setCJKmonofont{Consolas}}{\setCJKmonofont{Noto Sans Mono CJK SC}}
\IfFontExistsTF{Times New Roman}{\setmainfont{Times New Roman}}{\setmainfont{TeX Gyre Termes}}
\IfFontExistsTF{Arial}{\setsansfont{Arial}}{\setsansfont{TeX Gyre Heros}}

% ---------- Packages ----------
\usepackage{amsmath,amssymb}
\usepackage{booktabs}
\usepackage{colortbl}
\usepackage{array}
\usepackage{tabularx}
\usepackage{makecell}
\usepackage{pifont}
\usepackage{tikz}
\usepackage{hyperref}
\usetikzlibrary{arrows.meta,positioning,calc,fit,shapes.geometric}

\setlength{\emergencystretch}{2em}
\hfuzz=3pt
\setlength{\parskip}{2pt}
\graphicspath{{materials/figures/}{./}}
\AtBeginSection[]{}

% ---------- Local styles ----------
\definecolor{rvgreen}{RGB}{24,118,67}
\definecolor{rvred}{RGB}{164,42,42}
\definecolor{rvorange}{RGB}{181,103,22}
\definecolor{rvblue2}{RGB}{41,91,140}
\definecolor{rvgraybg}{RGB}{246,247,249}
\newcommand{\cmark}{\textcolor{rvgreen}{\ding{51}}}
\newcommand{\warnmark}{\textcolor{rvorange}{\ding{115}}}
\newcommand{\xmark}{\textcolor{rvred}{\ding{55}}}
\newcommand{\tagok}[1]{\textcolor{rvgreen}{\textbf{#1}}}
\newcommand{\tagwarn}[1]{\textcolor{rvorange}{\textbf{#1}}}
\newcommand{\tagno}[1]{\textcolor{rvred}{\textbf{#1}}}
\newcolumntype{Y}{>{\raggedright\arraybackslash}X}
\newcolumntype{L}[1]{>{\raggedright\arraybackslash}p{#1}}

\title[RV-MalTrace 论文进展]{RV-MalTrace 当前进展与论文推进计划}
\author[RV-MalTrace]{RV-MalTrace}
\institute[论文进展汇报]{论文进展汇报}
\date{2026年5月29日}
\setsupervisor{}
\setmajor{系统安全 / RISC-V 软硬件协同}

\begin{document}

% P1
\begin{frame}[plain]
  \begin{tikzpicture}[remember picture,overlay]
    \fill[accentcolor] (current page.north west) rectangle ([yshift=-0.18cm]current page.north east);
    \fill[accentcolor!8] ([yshift=-0.18cm]current page.north west) rectangle (current page.south east);
    \node[anchor=west,text=accentcolor,font=\sffamily\bfseries\fontsize{22pt}{28pt}\selectfont]
      at ([xshift=1.0cm,yshift=-1.55cm]current page.north west)
      {RV-MalTrace 当前进展};
    \node[anchor=west,text=inkcolor,font=\sffamily\bfseries\fontsize{17pt}{23pt}\selectfont]
      at ([xshift=1.0cm,yshift=-2.35cm]current page.north west)
      {35T 硬件追踪证据链与论文推进路线};
    \draw[accentcolor,line width=0.8pt]
      ([xshift=1.0cm,yshift=-3.05cm]current page.north west) --
      ([xshift=-1.0cm,yshift=-3.05cm]current page.north east);

    \node[anchor=west,align=left,text=inkcolor,font=\sffamily\fontsize{10.5pt}{15pt}\selectfont]
      at ([xshift=1.05cm,yshift=-4.25cm]current page.north west) {
        \textbf{主题：}RISC-V 硬件辅助行为追踪\\[4pt]
        \textbf{类型：}论文进展汇报\\[4pt]
        \textbf{核心问题：}当前能做到什么、边界在哪里、下一步做什么
      };

    \node[anchor=south east,align=right,text=textgray,font=\sffamily\fontsize{9pt}{12pt}\selectfont]
      at ([xshift=-1.0cm,yshift=0.75cm]current page.south east) {2026年5月29日};
  \end{tikzpicture}
\end{frame}

% P2
\begin{frame}
  \frametitle{汇报提纲}
  \vskip0.25cm
  {\footnotesize
  \begin{tabbing}
  \hspace{0.45cm}\=\hspace{0.72cm}\=\kill
  \textbf{\color{accentcolor}一}\>\textbf{研究定位}\,——\,把 trace tap 放进论文故事线\\[8pt]
  \textbf{\color{accentcolor}二}\>\textbf{实验数据与结果}\,——\,trace 格式样例、各流程能力、数据与 gate\\[8pt]
  \textbf{\color{accentcolor}三}\>\textbf{证据边界}\,——\,现在能说什么，哪些还不能说\\[8pt]
  \textbf{\color{accentcolor}四}\>\textbf{下一阶段}\,——\,从原型闭环推进到论文贡献
  \end{tabbing}
  }
  \vskip0.25cm
  \keybox{\textbf{本次报告的重点不是罗列文件}\,——\,而是把当前仓库整理成可审阅的科研进展、风险边界和下一步 gate。}
\end{frame}

\section{研究定位}

% P3
\begin{frame}
  \frametitle{Outline}
  \vskip0.4cm
  \tableofcontents[currentsection]
\end{frame}

% P4
\begin{frame}
  \frametitle{这不是一个单纯的 RTL demo}
  \vskip0.15cm

  当前仓库已经从“在 RISC-V/CVA6 或 LiteX/VexRiscv 上加一个 trace tap”推进到\alert{可复现实验证据链}：事件格式、仿真 golden、35T 板级运行、行为恢复、解释接口和 claim boundary 都有对应文件与检查脚本。

  \vskip0.22cm

  对论文进展报告而言，最稳妥的表达是：RV-MalTrace 现在可以支撑一个\alert{35T / LiteX / VexRiscv 硬件追踪辅助的恶意行为证据链原型}。这不是成熟检测器，也不是 CVA6 板级结论，但已经足够作为下一步论文路线的工程底座。

  \vskip0.2cm
  \keybox{\textbf{一句话进展：}当前原型已完成受控行为矩阵的板级 trace 证据闭环，并开始把真实恶意代码参考中的行为抽象迁移到安全可执行样例。}
\end{frame}

% P5
\begin{frame}
  \frametitle{当前技术链路（图 1）}
  \vskip0.05cm
  \begin{center}
  \begin{tikzpicture}[
    scale=0.94,transform shape,
    node distance=0.45cm and 0.75cm,
    box/.style={draw=lightline,fill=white,rounded corners=2pt,minimum height=0.72cm,align=center,font=\scriptsize,text width=2.30cm},
    wide/.style={draw=accentcolor!50,fill=accentcolor!7,rounded corners=2pt,minimum height=0.80cm,align=center,font=\scriptsize,text width=2.85cm},
    arrow/.style={-{Latex[length=2.0mm]},line width=0.55pt,color=accentcolor}
  ]
    \node[wide] (core) {RISC-V core\\CVA6 / VexRiscv};
    \node[box,right=of core] (tap) {commit / trap / syscall\\sideband tap};
    \node[box,right=of tap] (packet) {event packet\\formatter + filter};
    \node[box,right=of packet] (sink) {FIFO / BRAM / UART\\DROP accounting};

    \node[box,below=0.9cm of tap] (sim) {Vivado xsim\\golden JSONL};
    \node[box,right=of sim] (board) {Artix-7 35T\\512-record policy};
    \node[box,right=of board] (decode) {host decoder\\trace-code join};

    \node[wide,below=0.9cm of board] (sem) {behavior recovery\\fd/path, process tree, rule audit};
    \node[wide,right=of sem] (evid) {evidence package\\hash, boundary, checklist};

    \draw[arrow] (core) -- (tap);
    \draw[arrow] (tap) -- (packet);
    \draw[arrow] (packet) -- (sink);
    \draw[arrow] (tap) -- (sim);
    \draw[arrow] (sink) -- (board);
    \draw[arrow] (board) -- (decode);
    \draw[arrow] (decode) -- (sem);
    \draw[arrow] (sem) -- (evid);
  \end{tikzpicture}
  \figcap{图 1\;当前仓库覆盖从硬件事件采集到行为证据包的完整原型链路。}
  \end{center}
  \vskip0.05cm
  \small
  这条链路的关键点是\alert{低扰动事件采集}和\alert{离线证据解释}分离：硬件侧只保留提交事件与必要上下文，复杂语义在 host 侧恢复并受 gate 约束。
\end{frame}

\section{实验数据与结果}

% P6
\begin{frame}
  \frametitle{Outline}
  \vskip0.4cm
  \tableofcontents[currentsection]
\end{frame}

% P7
\begin{frame}
  \frametitle{仓库已经形成的四层能力（表 1）}
  \vskip0.1cm
  \begin{center}\scriptsize
  \renewcommand{\arraystretch}{1.22}
  \begin{tabularx}{0.96\linewidth}{@{}L{2.25cm}Y L{2.55cm}@{}}
    \toprule
    \textbf{层次} & \textbf{现在已经能做的事} & \textbf{主要证据} \\
    \midrule
    trace RTL / 格式 & 记录已提交、控制流、syscall、trap、context、\texttt{ARG\_MEM}、\texttt{DROP} 与 \texttt{MARKER}，并有过滤、压缩原型、DROP 计数。 & \texttt{trace\_format.md} \\
    \rowcolor{black!4}
    Vivado / CVA6 仿真 & trace unit、RVFI adapter、direct-core CVA6、full SoC probe、RV64GC microprobe 均有 PASS 记录。 & \texttt{sim\_results.md} \\
    35T 板级实验 & Artix-7 35T / LiteX / VexRiscv 下，13 个受控样例 full matrix 通过 512-record gate。 & 35T evidence bundle \\
    \rowcolor{black!4}
    行为解释与论文证据 & 支持 code map、trace-code join、fd/path、process tree、real-malware-derived lineage 和 claim boundary。 & paper evidence chain \\
    \bottomrule
  \end{tabularx}
  \end{center}

  \vskip0.12cm
  \small
  这四层已经让项目从“能跑一个样例”变成“能解释、能复现、能限制说法”的研究原型。
\end{frame}

% P7b
\begin{frame}
  \frametitle{阶段实验数据总览（表 2）}
  \vskip0.05cm
  \begin{center}\tiny
  \renewcommand{\arraystretch}{1.06}
  \begin{tabularx}{0.98\linewidth}{@{}L{2.15cm}L{2.15cm}L{3.65cm}Y@{}}
    \toprule
    \textbf{流程/阶段} & \textbf{当前程度} & \textbf{实验数据} & \textbf{结果解释} \\
    \midrule
    Trace / Vivado & \tagok{PASS} & 30 个 PASS 行；220 events；覆盖 syscall return、pointer、backpressure。 & packet 语义、DROP 计数和 CVA6 仿真入口可复现。 \\
    \rowcolor{black!4}
    35T 主矩阵 & \tagok{13/13 PASS} & 512 records；1,685 entry + 1,685 ret；130 marker；596 trap；max DROP 3.07\%。 & 受控 benign + synthetic malware-like gate 闭环。 \\
    语义闭环 & \tagok{case-study PASS} & fd/path 3/3 PASS，7 closed flows；process-chain 2 条边；function 6/6。 & 代表性解释路径可复述，不是完整语义恢复。 \\
    \rowcolor{black!4}
    对比流程 & \tagok{8 PASS / 2 deferred} & host/QEMU/strace/eBPF/QEMU-plugin/software/RVMT event-only 均有 13 样例证据。 & baseline 框架已具备，pointer/helper 待补。 \\
    真实恶意派生 & \tagok{6/6 PASS} & DarthRa/Mirai-derived；events 48--66；DROP 0；board 95.5--113.2 ms。 & 支持“参考行为可追踪”，不支持家族检测准确率。 \\
    \rowcolor{black!4}
    资源/非干扰 & \tagok{CHECKED} & LUT +48.04\%，FF +4.97\%，BRAM/DSP +0，slack 0.177 ns；6 项检查。 & 可引用实现成本和不反压原则，不能包装成性能提升。 \\
    \bottomrule
  \end{tabularx}
  \end{center}
\end{frame}

% P8
\begin{frame}
  \frametitle{Trace 事件模型已经稳定}
  \vskip0.1cm

  当前事件模型围绕\alert{已提交行为}设计，而不是完整 load/store 记录。syscall entry 从 U-mode ECALL 的 trap path 捕获，syscall return 从 SRET 返回 U-mode 捕获，DROP 作为一等事件进入 trace，避免在带宽不足时把 core 停住。

  \vskip0.22cm
  \begin{center}\scriptsize
  \setlength{\tabcolsep}{5pt}
  \begin{tabular}{@{}lll@{}}
    \toprule
    \textbf{事件族} & \textbf{代表事件} & \textbf{用于回答的问题} \\
    \midrule
    控制流 & \texttt{RETIRE, BRANCH, JUMP} & 执行路径、压缩指令长度、跳转目标 \\
    syscall & \texttt{SYSCALL\_ENTRY, SYSCALL\_RET} & 参数、返回值、持续周期、系统调用序列 \\
    异常与上下文 & \texttt{TRAP, CSR, SATP, PRIV} & trap cause、地址空间、特权级切换 \\
    资源边界 & \texttt{ARG\_MEM, DROP, MARKER} & 指针快照、溢出可见性、样例作用域 \\
    \bottomrule
  \end{tabular}
  \end{center}

  \vskip0.18cm
  \keybox{\textbf{设计取舍：}默认不追踪完整内存访问；需要语义时用 syscall-scoped \texttt{ARG\_MEM} 或 side-channel 进行有界补全。}
\end{frame}

% P8a
\begin{frame}
  \frametitle{Trace JSONL 的具体格式（示例 1）}
  \vskip0.05cm
  \small
  35T 板级 trace 经 host converter 后输出为\alert{JSON Lines}：每行一条事件记录。公共字段包括 \texttt{record\_index}、\texttt{cycle}、\texttt{evt}、\texttt{evt\_code}、\texttt{pc}、\texttt{instr}、\texttt{priv}、\texttt{raw\_header}、\texttt{raw\_words}；事件特有字段按需出现。

  \vskip0.12cm
  {\scriptsize\textbf{样例：}\texttt{benign/hello/board/trace-on/rep\_00/trace.jsonl} 的前三行解码结果。}
  \vskip0.08cm
  \begin{center}\tiny
  \renewcommand{\arraystretch}{1.10}
  \setlength{\tabcolsep}{3pt}
  \begin{tabularx}{0.98\linewidth}{@{}L{0.55cm}L{1.75cm}L{1.45cm}L{2.0cm}Y@{}}
    \toprule
    \textbf{\#} & \textbf{evt(code)} & \textbf{cycle} & \textbf{pc / priv} & \textbf{该行实际 payload} \\
    \midrule
    0 & \texttt{MARKER(12)} & 19,457,332 & \texttt{0x00000000 / --} & \texttt{value=0xb0002cab}; \texttt{raw\_header=0x0000000c}; \texttt{raw\_words}=16 个 32-bit word。 \\
    \rowcolor{black!4}
    1 & \texttt{SYSCALL\_ENTRY(4)} & 19,457,680 & \texttt{0x0101f554 / U} & \texttt{instr=0x00000073}; \texttt{syscall\_id=0x0000b380}; \texttt{a0=1}; \texttt{a1=0x9d1c3ad8}; \texttt{a7=0}。 \\
    2 & \texttt{SYSCALL\_RET(5)} & 19,464,490 & \texttt{0xc068d6f4 / S} & \texttt{duration=6810}; \texttt{target=0x0101f558}; \texttt{a0=0}; \texttt{a7=0x193}; \texttt{raw\_header=0x00000055}。 \\
    \bottomrule
  \end{tabularx}
  \end{center}

  \vskip0.06cm
  \keybox{\textbf{可检查点：}PPT 中的 trace 格式不是抽象引用；这里展示的是板级 artifact 中实际 JSONL 行的字段化内容。}
\end{frame}

% P8b
\begin{frame}
  \frametitle{16-word raw record 到 JSONL 的映射（示例 2）}
  \vskip0.05cm
  \small
  原始记录层保留为固定 16 个 32-bit word；JSONL 是在不丢弃 \texttt{raw\_words} 的前提下，把常用语义字段展开出来。下面同样来自 \texttt{hello} 的第 1 条 \texttt{SYSCALL\_ENTRY}。

  \vskip0.12cm
  \begin{center}\tiny
  \renewcommand{\arraystretch}{1.14}
  \setlength{\tabcolsep}{3pt}
  \begin{tabularx}{0.96\linewidth}{@{}L{1.65cm}L{3.65cm}Y@{}}
    \toprule
    \textbf{raw word} & \textbf{实际值} & \textbf{JSONL 字段含义} \\
    \midrule
    \texttt{w0} & \texttt{0x00000004} & \texttt{raw\_header}; 低位解码为 \texttt{evt\_code=4}，即 \texttt{SYSCALL\_ENTRY}。 \\
    \rowcolor{black!4}
    \texttt{w1} & \texttt{0x0128e690} & \texttt{cycle=19457680}，用于与 return、DROP、marker 做时序关联。 \\
    \texttt{w2, w3} & \texttt{0x0101f554}, \texttt{0x00000073} & \texttt{pc=0x000000000101f554}; \texttt{instr=ECALL}。 \\
    \rowcolor{black!4}
    \texttt{w6} & \texttt{0x0000b380} & \texttt{syscall\_id}，用于 entry/return 配对。 \\
    \texttt{w8--w15} & \makecell[l]{\texttt{a0=0x00000001}, \texttt{a1=0x9d1c3ad8}\\\texttt{a2=0x0}, \texttt{a3=0x1}\\\texttt{a4=0x000f423f}, \texttt{a5=0x0}\\\texttt{a6=0x0}, \texttt{a7=0x0}} & syscall 参数寄存器快照；后续 fd/path、规则审计、baseline 对齐均基于这些字段。 \\
    \bottomrule
  \end{tabularx}
  \end{center}

  \vskip0.08cm
  {\footnotesize\textbf{格式结论：}trace 文件采用“语义 JSONL + 原始 16-word 记录并存”的格式；既能做实验统计，也能回到 raw word 层复核转换。}
\end{frame}

% P9
\begin{frame}
  \frametitle{Vivado 与 CVA6 仿真门禁（表 3）}
  \vskip0.1cm
  \begin{center}\scriptsize
  \renewcommand{\arraystretch}{1.20}
  \begin{tabularx}{0.96\linewidth}{@{}L{3.15cm}L{1.35cm}Y@{}}
    \toprule
    \textbf{门禁} & \textbf{状态} & \textbf{说明} \\
    \midrule
    trace unit + RVFI adapter & \tagok{PASS} & 覆盖 syscall return、pointer string、pointer guardrails、backpressure、filter、board minimal。 \\
    \rowcolor{black!4}
    CVA6 direct-core matrix & \tagok{PASS} & smoke、branch、jump、ecall、illegal trap、ebreak 等 trace-on/no-trace 终态一致。 \\
    full SoC harness probes & \tagok{PASS} & breakpoint smoke、UART/MMIO store path、normal tohost/MMIO completion 已记录。 \\
    \rowcolor{black!4}
    RV64GC microprobe & \tagok{PASS} & I/M/A/F/D/C 每类至少一条 committed instruction probe 通过。 \\
    Linux / arbitrary program & \tagwarn{未声称} & 当前不是 Linux 全系统任意程序正确性结论，也不是 CVA6 物理板结论。 \\
    \bottomrule
  \end{tabularx}
  \end{center}

  \vskip0.1cm
  \small
  这部分适合在报告中说明：CVA6 仿真方向已有可复现基础，但论文级板上 Linux 语义仍要单独 gate。
\end{frame}

% P9b
\begin{frame}
  \frametitle{仿真流程实验数据（表 4）}
  \vskip0.05cm
  \begin{center}\scriptsize
  \renewcommand{\arraystretch}{1.14}
  \begin{tabularx}{0.97\linewidth}{@{}L{2.9cm}L{1.65cm}L{3.55cm}Y@{}}
    \toprule
    \textbf{子流程} & \textbf{行数/状态} & \textbf{代表数据} & \textbf{实验结果} \\
    \midrule
    trace unit semantics & 12 / \tagok{PASS} & \texttt{pointer\_string}: 13 events / 5 ARG\_MEM；pointer guardrails: 46 events / 14 ARG\_MEM。 & syscall pointer snapshot 的 synthetic 路径和保护条件通过。 \\
    \rowcolor{black!4}
    overflow / filter & 2 / \tagok{PASS} & \texttt{backpressure}: 30 events，7 DROP records，18 dropped events；\texttt{filter}: 仅保留 1 branch。 & 溢出可见且不反压 core，过滤逻辑可控。 \\
    RVFI adapter & 1 / \tagok{PASS} & 15 events，覆盖 dual commit、U-mode entry/return、non-ECALL trap、compressed control flow。 & CVA6 RVFI 到 trace schema 的适配已单测通过。 \\
    \rowcolor{black!4}
    CVA6 direct-core & 6 / \tagok{PASS} & smoke、branch、jump、ecall、illegal、ebreak；trace-on/no-trace final PASS 对齐。 & 真实 CVA6 committed RVFI 事件可进入同一 checker。 \\
    full SoC + RV64GC & 9 / \tagok{PASS} & full SoC smoke/store/tohost + I/M/A/F/D/C microprobe；30 行总计 220 events。 & full SoC 最小探针与 RV64GC 扩展族最小 retire gate 已覆盖。 \\
    \bottomrule
  \end{tabularx}
  \end{center}
\end{frame}

% P10
\begin{frame}
  \frametitle{35T 主矩阵已经闭环（表 5）}
  \vskip0.1cm
  \begin{center}\scriptsize
  \renewcommand{\arraystretch}{1.18}
  \begin{tabularx}{0.96\linewidth}{@{}L{2.75cm}L{2.65cm}Y@{}}
    \toprule
    \textbf{指标} & \textbf{结果} & \textbf{解释} \\
    \midrule
    run id & \makecell[l]{\texttt{35t-smallcap-r512}\\\texttt{full matrix}\\\texttt{20260521}} & 主 evidence bundle，用于严格 trace gate。 \\
    \rowcolor{black!4}
    样例矩阵 & \tagok{13/13 PASS} & benign + synthetic malware-like，全矩阵 gate 通过。 \\
    trace 容量 & 512 records & 并非通过扩大到 1024 records 解决；当前按 small-capacity policy 分样例配置。 \\
    \rowcolor{black!4}
    trace health & \tagok{UNKNOWN = 0} & marker scope、runtime process attribution、DROP/capacity 均进入 gate；corrupt event 也为 0。 \\
    语义补全 & \tagok{selected PASS} & targeted validation 关闭 fd/path 与 process-tree 代表性解释路径。 \\
    \bottomrule
  \end{tabularx}
  \end{center}

  \vskip0.13cm
  \small
  这里最有价值的是“\alert{35T 小容量约束下仍可建立证据链}”，不是成熟恶意软件检测能力。
\end{frame}

% P10b
\begin{frame}
  \frametitle{35T 样例矩阵实验数据（表 6）}
  \vskip0.05cm
  \begin{center}\tiny
  \renewcommand{\arraystretch}{1.05}
  \begin{tabularx}{0.97\linewidth}{@{}L{2.7cm}L{1.35cm}L{2.25cm}L{2.0cm}Y@{}}
    \toprule
    \textbf{样例组} & \textbf{样例数} & \textbf{事件计数} & \textbf{DROP} & \textbf{结果} \\
    \midrule
    benign 基线 & 5 & 545 entry / 545 ret / 50 marker & max 0 & \texttt{hello, ls, cat, cp, sha256sum} 全部 PASS；\texttt{ls} 的 file-scan overlap 被标注为 benign overlap。 \\
    \rowcolor{black!4}
    synthetic malware-like & 8 & 1,140 entry / 1,140 ret / 80 marker / 596 trap & max 3.07\% & 8/8 PASS；\texttt{illegal\_trap} 使用 \texttt{p0c}，其余使用 \texttt{p0a}。 \\
    capacity-sensitive case & 1 & \texttt{process\_chain}: 380 entry / 380 ret & 0 & 不再是 512-record blocker；每个 rep 154 events，无 cap hit。 \\
    \rowcolor{black!4}
    trap-heavy case & 1 & \texttt{illegal\_trap}: 90 entry / 90 ret / 596 trap & 3.07\% & trap profile 捕获 illegal-instruction 行为，仍低于 5\% gate。 \\
    aggregate gate & 13 & 3,370 syscall records + 130 marker + 596 trap & max 3.07\% & 13/13 strict sample gate PASS，UNKNOWN/corrupt 为 0。 \\
    \bottomrule
  \end{tabularx}
  \end{center}
\end{frame}

% P10c
\begin{frame}
  \frametitle{\texttt{illegal\_trap} 的板级结果样例（示例 3）}
  \vskip0.04cm
  \begin{columns}[T,onlytextwidth]
    \column{0.55\textwidth}
    {\scriptsize\textbf{trace.jsonl 片段}}
    \vskip0.05cm
    \begin{center}\tiny
    \renewcommand{\arraystretch}{1.06}
    \setlength{\tabcolsep}{2.5pt}
    \begin{tabularx}{0.99\linewidth}{@{}L{0.6cm}L{1.75cm}Y@{}}
      \toprule
      \textbf{\#} & \textbf{事件} & \textbf{实际字段} \\
      \midrule
      0 & \texttt{MARKER} & \texttt{value=0xb0000f79}; sample begin。 \\
      \rowcolor{black!4}
      1 & \texttt{SYSCALL\_ENTRY} & \texttt{pc=0x0101f554}; \texttt{syscall\_id=0x0000e101}; \texttt{a0=1}。 \\
      2 & \texttt{SYSCALL\_RET} & \texttt{duration=6819}; \texttt{target=0x0101f558}; \texttt{a0=0}。 \\
      \rowcolor{black!4}
      5 & \texttt{TRAP} & \texttt{cause=0x9}; \texttt{tval=0x73}; \texttt{pc=0xc000788c}; \texttt{priv=S}。 \\
      6 & \texttt{TRAP} & \texttt{cause=0x2}; \texttt{tval=0x60059593}; supervisor trap。 \\
      \rowcolor{black!4}
      13 & \texttt{TRAP} & \texttt{cause=0xc}; \texttt{pc=0x0001059c}; \texttt{priv=U}。 \\
      \bottomrule
    \end{tabularx}
    \end{center}

    \column{0.45\textwidth}
    {\scriptsize\textbf{trace-code join 摘要}}
    \vskip0.05cm
    \begin{center}\tiny
    \renewcommand{\arraystretch}{1.08}
    \setlength{\tabcolsep}{3pt}
    \begin{tabularx}{0.98\linewidth}{@{}L{2.2cm}Y@{}}
      \toprule
      \textbf{字段} & \textbf{实际结果} \\
      \midrule
      events & 158 \\
      marker scope & \texttt{PASS}; begin \texttt{0xb0000f79}, end \texttt{0xe0000f79} \\
      runtime path & \texttt{/usr/bin/illegal\_trap}; runtime map \texttt{PASS} \\
      PC owner & kernel 121, target 27, unknown 10 \\
      callsite kind & illegal-instruction site 2; syscall site 14 \\
      DROP & max 3.07\%, below 5\% gate \\
      \bottomrule
    \end{tabularx}
    \end{center}
  \end{columns}

  \vskip0.05cm
  \keybox{\textbf{判定顺序：}先看到 trap-heavy trace、marker 闭合、runtime/process 归因与 code-map 结果，再给出结论：该样例在 512-record policy 下通过板级 gate。}
\end{frame}

% P11
\begin{frame}
  \frametitle{真实恶意代码派生行为有六行证据（表 7）}
  \vskip0.05cm
  \begin{center}\scriptsize
  \renewcommand{\arraystretch}{1.12}
  \begin{tabularx}{0.98\linewidth}{@{}L{3.45cm}L{2.35cm}L{1.05cm}Y@{}}
    \toprule
    \textbf{行为行} & \textbf{来源} & \textbf{状态} & \textbf{边界} \\
    \midrule
    ELF header probe & DarthRa reference & \tagok{PASS} & 真实恶意代码派生行为，不是 payload 等价执行。 \\
    \rowcolor{black!4}
    rootkit device probe & DarthRa reference & \tagok{PASS} & 安全控制下的行为验证。 \\
    virus fixture walk & DarthRa reference & \tagok{PASS} & 文件遍历类行为抽象。 \\
    \rowcolor{black!4}
    proc scan & Mirai reference & \tagok{PASS} & 非网络版本，排除外联。 \\
    watchdog probe & Mirai reference & \tagok{PASS} & 非网络行为抽象。 \\
    \rowcolor{black!4}
    encoded table & Mirai reference & \tagok{PASS} & 字符串/表行为抽象。 \\
    \bottomrule
  \end{tabularx}
  \end{center}

  \vskip0.1cm
  \keybox{\textbf{可讲法：}这些结果支持“选取自真实恶意代码参考的行为可以被 35T 原型 trace、验证和规则审计”，不支持“真实家族检测准确率”。}
\end{frame}

% P11b
\begin{frame}
  \frametitle{真实恶意派生流程的时序与 trace 数据（表 8）}
  \vskip0.05cm
  \begin{center}\scriptsize
  \renewcommand{\arraystretch}{1.10}
  \setlength{\tabcolsep}{3.5pt}
  \begin{tabularx}{0.98\linewidth}{@{}L{3.25cm}rrr r r L{1.25cm}@{}}
    \toprule
    \textbf{样例} & \textbf{Host ms} & \textbf{QEMU ms} & \textbf{Board ms} & \textbf{Events} & \textbf{DROP} & \textbf{状态} \\
    \midrule
    darthra\_elf\_header\_probe & 17.3 & 27.7 & 95.5 & 54 & 0 & PASS \\
    \rowcolor{black!4}
    darthra\_rootkit\_device\_probe & 17.7 & 28.8 & 109.7 & 52 & 0 & PASS \\
    darthra\_virus\_fixture\_walk & 29.4 & 73.5 & 113.2 & 66 & 0 & PASS \\
    \rowcolor{black!4}
    mirai\_proc\_scan & 12.4 & 20.8 & 105.2 & 56 & 0 & PASS \\
    mirai\_watchdog\_probe & 12.2 & 19.1 & 97.8 & 48 & 0 & PASS \\
    \rowcolor{black!4}
    mirai\_encoded\_table & 18.5 & 44.1 & 100.5 & 50 & 0 & PASS \\
    \bottomrule
  \end{tabularx}
  \end{center}

  \vskip0.08cm
  \small
  这些数据用于说明同一组安全控制行为在 host、QEMU 和 35T trace-on 三条路径下均有对应记录；比值是描述性结果，不是性能等价或加速结论。
\end{frame}

% P12
\begin{frame}
  \frametitle{解释接口让结果可复述}
  \vskip0.05cm
  \begin{columns}[T,onlytextwidth]
    \column{0.52\textwidth}
    \small
    现在的结果不是只能看原始 UART 或 JSONL。仓库已经提供终端解释接口，把 trace health、gate 状态、恢复出的行为、可疑 cue、证据来源和 non-claim 一起打印。

    \vskip0.2cm
    \begin{center}
    \scriptsize
    \begin{tabular}{@{}ll@{}}
      \toprule
      \textbf{命令} & \textbf{用途} \\
      \midrule
      \texttt{rvmt explain:35t} & 单样例解释 \\
      \texttt{--flow} & run-level 过程视图 \\
      \texttt{--format markdown} & 生成可提交材料 \\
      \bottomrule
    \end{tabular}
    \end{center}

    \column{0.48\textwidth}
    \small
    目前解释层已经覆盖三类代表性语义：

    \vskip0.2cm
    \begin{enumerate}\setlength\itemsep{0.45em}
      \item \textbf{code map}\,——\,PC 与 ELF symbol / syscall site 关联；
      \item \textbf{fd/path}\,——\,代表性 \texttt{openat -> fd -> read/write/close} 闭环；
      \item \textbf{process tree}\,——\,\texttt{clone} 返回 child PID 与 \texttt{wait} PID 匹配。
    \end{enumerate}
  \end{columns}
\end{frame}

% P12b
\begin{frame}
  \frametitle{语义恢复流程的实验结果（表 9）}
  \vskip0.05cm
  \begin{center}\scriptsize
  \renewcommand{\arraystretch}{1.13}
  \begin{tabularx}{0.97\linewidth}{@{}L{2.75cm}L{3.45cm}L{2.15cm}Y@{}}
    \toprule
    \textbf{流程} & \textbf{实验数据} & \textbf{状态} & \textbf{结果边界} \\
    \midrule
    fd/path flow & 3/3 required samples PASS；\texttt{file\_scan}=1 closed flow，\texttt{batch}=4，\texttt{self\_copy}=2；unresolved fd = 0。 & \tagok{PASS} & path string 来源是 board syscall side-channel，不是硬件 pointer snapshot。 \\
    \rowcolor{black!4}
    process tree & \texttt{process\_chain}: 10 candidates，5 PASS candidates；2 条 clone-return/waitid 边；2 个 \texttt{/bin/true} exec path。 & \tagok{PASS} & parent PID 仍为 \texttt{target\_parent\_unresolved}，不声称完整 OS process graph。 \\
    function attribution & 6/6 case-study samples available；ELF symbol table；约 1,099--1,103 functions/sample；119 target-attributed events。 & \tagok{PASS} & function-level 是 symbol/range 证据，不等于 source-line。 \\
    \rowcolor{black!4}
    source-line attribution & 6 个 case-study samples 均记录 \texttt{source\_line\_count=0}；source line basis unavailable。 & \tagwarn{PARTIAL} & 下一步需要保留 DWARF/source-location records。 \\
    \bottomrule
  \end{tabularx}
  \end{center}
\end{frame}

% P12c
\begin{frame}
  \frametitle{\texttt{file\_scan} 的 fd/path 输出样例（示例 4）}
  \vskip0.05cm
  \small
  下面不是“fd/path PASS”的文字概括，而是 canonical summary 选中的 board syscall side-channel 结果。它展示了路径、fd 生成、目录读取和 close 如何闭合成一个 flow。

  \vskip0.10cm
  \begin{center}\tiny
  \renewcommand{\arraystretch}{1.12}
  \setlength{\tabcolsep}{3pt}
  \begin{tabularx}{0.98\linewidth}{@{}L{0.7cm}L{1.25cm}L{0.95cm}L{1.45cm}Y@{}}
    \toprule
    \textbf{seq} & \textbf{name} & \textbf{fd} & \textbf{return} & \textbf{结果字段} \\
    \midrule
    12 & \texttt{openat} & 3 & \texttt{0x3} & \texttt{path=experiments/linux\_behavior/malware\_like/fixtures/scan\_root}; \texttt{path\_source=board\_syscall\_side\_channel}。 \\
    \rowcolor{black!4}
    13 & \texttt{getdents64} & 3 & \texttt{0x70} & 目录项读取返回非零长度。 \\
    14 & \texttt{getdents64} & 3 & \texttt{0x0} & 第二次读取到目录末尾。 \\
    \rowcolor{black!4}
    15 & \texttt{close} & 3 & \texttt{0x0} & fd 生命周期闭合；\texttt{fd\_generation=1}。 \\
    \bottomrule
  \end{tabularx}
  \end{center}

  \vskip0.08cm
  \keybox{\textbf{实验结果：}\texttt{openat(path) -> fd=3 -> getdents64 -> getdents64 -> close} 被恢复为 \texttt{status=closed}；\texttt{unresolved\_fds=0}，因此 fd/path flow 在该样例上判定为 PASS。}
\end{frame}

% P12d
\begin{frame}
  \frametitle{\texttt{process\_chain} 的 process-tree 输出样例（示例 5）}
  \vskip0.05cm
  \small
  process-tree 结果来自同一类 canonical summary：选取最强 side-channel candidate，显式输出 clone 返回值、wait 目标 PID、exec path 和最终边。

  \vskip0.10cm
  \begin{center}\tiny
  \renewcommand{\arraystretch}{1.12}
  \setlength{\tabcolsep}{3pt}
  \begin{tabular}{@{}L{1.8cm}L{2.35cm}L{7.1cm}@{}}
    \toprule
    \textbf{证据项} & \textbf{实际输出} & \textbf{解释} \\
    \midrule
    observed counts & \makecell[l]{clone/fork = 2\\execve = 4\\wait = 2} & 目标样例中观察到两条进程创建链。 \\
    \rowcolor{black!4}
    clone seq 12 & \makecell[l]{return = \texttt{0xcb}\\child PID = 203} & parent-side clone return 被解释为子进程 PID。 \\
    wait seq 13 & \makecell[l]{wait PID = 203\\return = \texttt{0x0}} & wait 目标 PID 与 clone 返回 PID 匹配。 \\
    \rowcolor{black!4}
    clone seq 14 & \makecell[l]{return = \texttt{0xcc}\\child PID = 204} & 第二条子进程边。 \\
    wait seq 15 & \makecell[l]{wait PID = 204\\return = \texttt{0x0}} & 第二条 wait 匹配闭合。 \\
    \rowcolor{black!4}
    exec path & \makecell[l]{\texttt{/bin/true}\\source=board\\syscall side-channel} & exec path 来源明确，不是文件名猜测。 \\
    edges & \makecell[l]{parent unres. \(\rightarrow\) 203\\parent unres. \(\rightarrow\) 204\\confidence: strong} & 父进程身份按边界保守输出为 unresolved，不强行补全 parent PID。 \\
    \bottomrule
  \end{tabular}
  \end{center}

  \vskip0.06cm
  {\footnotesize\textbf{实验结果：}summary status 为 PASS，partial edges 与 unclosed edges 均为空；但 parent PID 仍为 unresolved，不声称完整 OS process graph。}
\end{frame}

\begin{frame}
  \frametitle{对比实验流程当前数据（表 10）}
  \vskip0.05cm
  \begin{center}\scriptsize
  \renewcommand{\arraystretch}{1.13}
  \begin{tabularx}{0.97\linewidth}{@{}L{2.9cm}L{1.35cm}L{2.6cm}Y@{}}
    \toprule
    \textbf{baseline / ablation} & \textbf{状态} & \textbf{样例覆盖} & \textbf{当前结果} \\
    \midrule
    host native / host strace & \tagok{PASS} & 13/13 & host strace/native median = 2.60；用于说明软件 syscall 路径扰动。 \\
    \rowcolor{black!4}
    QEMU native / QEMU strace & \tagok{PASS} & 13/13 & qemu strace/native median = 1.84；提供模拟器基线。 \\
    eBPF-only & \tagok{PASS} & 13/13 & bpftrace raw syscall tracepoint，comm-filtered host binaries。 \\
    \rowcolor{black!4}
    QEMU plugin & \tagok{PASS} & 13/13 & QEMU user-mode TCG plugin syscall-count baseline；不是硬件 trace。 \\
    software instrumentation & \tagok{PASS} & 13/13 & \texttt{-finstrument-functions} host binary；不提供 syscall 参数恢复。 \\
    \rowcolor{black!4}
    RVMT event-only & \tagok{PASS} & 13/13 & trace bytes/syscall median = 672.9；board trace on/off median = 0.566，仅作描述。 \\
    pointer snapshot / helper & \tagwarn{DEFERRED} & 0/13 & pointer snapshot 与 helper/ebPF companion 仍是后续增强路径。 \\
    \bottomrule
  \end{tabularx}
  \end{center}
\end{frame}

\section{证据边界}

% P13
\begin{frame}
  \frametitle{Outline}
  \vskip0.4cm
  \tableofcontents[currentsection]
\end{frame}

% P14
\begin{frame}
  \frametitle{现在可以稳妥陈述的结论（表 11）}
  \vskip0.1cm
  \begin{center}\scriptsize
  \renewcommand{\arraystretch}{1.20}
  \begin{tabularx}{0.96\linewidth}{@{}L{3.5cm}Y L{2.25cm}@{}}
    \toprule
    \textbf{可陈述结论} & \textbf{精确含义} & \textbf{证据位置} \\
    \midrule
    35T 原型闭环 & Artix-7 35T / LiteX / VexRiscv 上，受控样例的硬件 trace、运行时归因、规则审计已闭环。 & application closure \\
    \rowcolor{black!4}
    主矩阵 gate 通过 & 512-record policy 下，13 个样例 13/13 PASS，UNKNOWN/corrupt 为 0。 & full matrix bundle \\
    代表性语义闭环 & fd/path 与 process-tree 代表路径已在 targeted validation 中 PASS。 & semantic summaries \\
    \rowcolor{black!4}
    真实恶意派生行为可追踪 & DarthRa/Mirai reference 的 6 行安全控制行为验证 PASS。 & lineage / baseline \\
    证据链纪律可用 & hash manifest、claim boundary、review checklist 已形成 CCF-A-style discipline。 & strong chain package \\
    \bottomrule
  \end{tabularx}
  \end{center}

  \vskip0.1cm
  \small
  这些说法的共同特点是：\alert{限定平台、限定样例、限定安全控制、限定证据类型}。
\end{frame}

% P15
\begin{frame}
  \frametitle{仍然不能越界的说法（表 12）}
  \vskip0.1cm
  \begin{center}\scriptsize
  \renewcommand{\arraystretch}{1.20}
  \begin{tabularx}{0.96\linewidth}{@{}L{3.1cm}L{1.15cm}Y@{}}
    \toprule
    \textbf{不能声称} & \textbf{状态} & \textbf{原因} \\
    \midrule
    CVA6 物理板验证完成 & \tagno{NO} & 当前 35T 结果来自 LiteX/VexRiscv，CVA6 主要是仿真和综合/资源路径。 \\
    \rowcolor{black!4}
    真实恶意软件检测准确率 & \tagno{NO} & 真实恶意派生行为不是家族分类、IOC/TTP 覆盖或检测质量评估。 \\
    完整语义恢复 & \tagno{NO} & fd/path 与 process-tree 是代表性闭环；source-line、全局进程所有权、完整内存语义仍缺。 \\
    \rowcolor{black!4}
    成熟检测器或产品级系统 & \tagno{NO} & 当前是 evidence-chain prototype，规则审计是应用展示，不是成熟 detector。 \\
    单 trace side-channel 全门禁 PASS & \tagno{NO} & targeted validation 是 dual-channel；side-channel 不是 strict full-matrix trace gate。 \\
    \bottomrule
  \end{tabularx}
  \end{center}

  \vskip0.1cm
  \keybox{\textbf{建议表述：}先讲已闭环证据，再主动讲 non-claims。这样能清楚呈现项目是“有边界的进展”，不是夸大。}
\end{frame}

% P16
\begin{frame}
  \frametitle{资源与非干扰证据的边界（表 13）}
  \vskip0.1cm
  \begin{center}\scriptsize
  \renewcommand{\arraystretch}{1.16}
  \begin{tabular}{@{}lrrr@{}}
    \toprule
    \textbf{指标} & \textbf{Baseline} & \textbf{Trace-enabled} & \textbf{Delta} \\
    \midrule
    LUT & 84,923 & 125,717 & +40,794 (+48.04\%) \\
    \rowcolor{black!4}
    FF & 56,491 & 59,301 & +2,810 (+4.97\%) \\
    BRAM18 equiv & 108 & 108 & +0 \\
    \rowcolor{black!4}
    DSP & 27 & 27 & +0 \\
    Slack & 0.177 ns & 0.177 ns & 0.000 ns \\
    \bottomrule
  \end{tabular}
  \end{center}

  \vskip0.18cm
  \small
  非干扰 gate 已检查 sideband-only、registered snapshot、DROP-not-stall、direct-core trace/no-trace final PASS 等条件。资源结果可以引用为\alert{当前 routed report 的 trace-enabled delta}。

  \vskip0.18cm
  \keybox{\textbf{边界：}这些数据不能被包装成 IPC/Fmax 提升或真实 workload overhead 结论；它们说明的是当前实现成本和非反压原则。}
\end{frame}

\section{下一阶段}

% P17
\begin{frame}
  \frametitle{Outline}
  \vskip0.4cm
  \tableofcontents[currentsection]
\end{frame}

% P18
\begin{frame}
  \frametitle{论文主线要从 prototype 走向 contribution}
  \vskip0.1cm

  如果目标只是“硬件能采 syscall event”，研究贡献会偏工程。更有论文竞争力的主线应当是：\alert{RISC-V 低扰动硬件辅助 syscall 语义追踪}，其中硬件 trace 提供 committed、低可见度的行为事实，离线语义层恢复 fd/path/process/memory-permission 等行为图。

  \vskip0.2cm

  下一阶段的中心问题不是再堆更多样例，而是回答：仅靠寄存器与上下文能恢复到哪里；哪些 pointer 参数必须补全；硬件 pointer snapshot、kernel helper/eBPF fallback 与 baseline 方法的边界分别是什么。

  \vskip0.2cm
  \keybox{\textbf{论文故事线：}从“35T 原型可行”推进到“RISC-V hardware-assisted semantic behavior tracing with bounded pointer reconstruction and evasion-aware evaluation”。}
\end{frame}

% P19
\begin{frame}
  \frametitle{下一步一：把 pointer 语义做成核心创新}
  \vskip0.1cm

  许多关键 syscall 的行为语义不在寄存器数值本身，而在 pointer 指向的用户内存中。当前 \texttt{ARG\_MEM} 已有 synthetic openat pathname 与 guardrails 回归，下一步要把它推进到 Linux/CVA6 信号接入或明确 fallback。

  \vskip0.1cm
  {\small
  \[
  \begin{aligned}
    &\mathrm{SYSCALL\_ENTRY}(a_7, a_0,\ldots,a_5)
      + \mathrm{ARG\_MEM}(ptr, bytes)\\
    &\qquad\Rightarrow\;
      \mathrm{openat}(\mathrm{path}),\;
      \mathrm{execve}(\mathrm{argv}),\;
      \mathrm{connect}(\mathrm{sockaddr})
  \end{aligned}
  \]}
  \vskip0.05cm

  \begin{center}\scriptsize
  \begin{tabularx}{0.92\linewidth}{@{}L{2.5cm}Y L{2.4cm}@{}}
    \toprule
    \textbf{路径} & \textbf{要验证的点} & \textbf{风险} \\
    \midrule
    hardware snapshot & S-mode copy-from-user load 能否稳定捕获字符串/结构体片段。 & LSU 信号、跨页、带宽 \\
    \rowcolor{black!4}
    kernel helper/eBPF & 只补 pid/path/string，不替代硬件 trace。 & 威胁模型收窄 \\
    offline reconstruction & fd/path/process/memory behavior graph 形成可审计输出。 & 完整性边界 \\
    \bottomrule
  \end{tabularx}
  \end{center}
\end{frame}

% P20
\begin{frame}
  \frametitle{下一步二：把对比实验补齐（表 14）}
  \vskip0.1cm
  \begin{center}\scriptsize
  \renewcommand{\arraystretch}{1.18}
  \begin{tabularx}{0.96\linewidth}{@{}L{2.65cm}Y L{2.3cm}@{}}
    \toprule
    \textbf{baseline} & \textbf{比较目标} & \textbf{关键指标} \\
    \midrule
    \texttt{strace/ptrace} & 最直接的软件 syscall ground truth 与 anti-debug 可见性。 & overhead, detectability \\
    \rowcolor{black!4}
    eBPF-only / helper & OS helper 路线的语义能力与扰动。 & semantic accuracy, overhead \\
    QEMU / plugin & 模拟器路径的可控性与真实执行差距。 & timing, coverage \\
    \rowcolor{black!4}
    software instrumentation & 编译期或二进制插桩与硬件旁路采集的对比。 & perturbation, completeness \\
    RV-MalTrace ablation & event-only、pointer snapshot、helper 三种模式拆开比较。 & accuracy, resource, drop \\
    \bottomrule
  \end{tabularx}
  \end{center}

  \vskip0.1cm
  \small
  对系统安全论文而言，实验不能只证明“能 trace”，还要证明\alert{为什么硬件辅助比常规软件方法有价值}。
\end{frame}

% P21
\begin{frame}
  \frametitle{未来两周优先级（表 15）}
  \vskip0.1cm
  \begin{center}\scriptsize
  \renewcommand{\arraystretch}{1.18}
  \begin{tabularx}{0.96\linewidth}{@{}L{2.6cm}Y L{2.05cm}@{}}
    \toprule
    \textbf{任务} & \textbf{交付物} & \textbf{判断标准} \\
    \midrule
    source-line attribution & 在 code map 中保留 DWARF/source line，并 join 到代表 syscall/trap site。 & source summary 从 PARTIAL 推进 \\
    \rowcolor{black!4}
    Linux semantic gate & 选择 2--3 个 open/exec/mmap/process workload，固定 runbook 与 ground truth。 & path/process 语义可复查 \\
    pointer snapshot 决策 & CVA6 LSU 信号 map + synthetic-to-Linux feasibility note。 & hardware/fallback 路线明确 \\
    \rowcolor{black!4}
    baseline plan & \texttt{strace}, QEMU, eBPF/helper 的可运行脚本和指标表。 & 同一样例可横向比较 \\
    paper claim table & 将 allowed claims / forbidden claims 变成论文式 contribution boundary。 & 可直接审阅 \\
    \bottomrule
  \end{tabularx}
  \end{center}

  \vskip0.08cm
  \keybox{\textbf{优先级原则：}先补“论文里一定会被问”的证据缺口，再扩样例规模。}
\end{frame}

% P21b
\begin{frame}
  \frametitle{未闭合流程的实验状态（表 16）}
  \vskip0.05cm
  \begin{center}\scriptsize
  \renewcommand{\arraystretch}{1.13}
  \begin{tabularx}{0.97\linewidth}{@{}L{2.9cm}L{2.7cm}L{2.35cm}Y@{}}
    \toprule
    \textbf{流程} & \textbf{当前数据} & \textbf{状态} & \textbf{下一步 gate} \\
    \midrule
    source-line attribution & 6 个 case-study samples：source line count = 0；function-level 6/6。 & \tagwarn{PARTIAL} & 保留 DWARF，输出 source-location records。 \\
    \rowcolor{black!4}
    hardware pointer snapshot & pointer snapshot samples\_with\_evidence = 0/13；synthetic \texttt{ARG\_MEM} 仿真 PASS。 & \tagwarn{DEFERRED} & CVA6 LSU map + Linux copy-from-user workload。 \\
    helper/eBPF companion & helper/eBPF companion evidence = 0/13；eBPF-only baseline 13/13 PASS。 & \tagwarn{DEFERRED} & 明确 trusted-kernel 边界，只补语义。 \\
    \rowcolor{black!4}
    side-channel single-run gate & sample status 13/13 PASS；strict sample gate 9/13 PASS。 & \tagwarn{not claimed} & 不能写成 single-trace all-gates PASS。 \\
    fuzz/stress validation & fuzz\_cf、fuzz\_trap、fuzz\_syscall、fuzz\_context、fuzz\_overflow：5 TODO。 & \tagwarn{TODO} & deterministic seeds + invariant checker。 \\
    \bottomrule
  \end{tabularx}
  \end{center}
\end{frame}

% P22
\begin{frame}
  \frametitle{需要讨论的三个决策}
  \vskip0.15cm

  接下来需要尽早确定方向，否则后续实验容易发散：

  \vskip0.2cm
  \begin{enumerate}\setlength\itemsep{0.55em}
    \item \textbf{论文定位}\,——\,主打 malware analysis，还是更稳地表述为 RISC-V semantic behavior tracing，并把 malware 作为应用场景；
    \item \textbf{语义补全路线}\,——\,优先冲 hardware pointer snapshot，还是 hardware trace + trusted helper 双路径并行；
    \item \textbf{目标强度}\,——\,当前 35T evidence chain 是否作为阶段性小论文/技术报告，还是继续补 CVA6/Linux/对比实验后冲更高层级会议。
  \end{enumerate}

  \vskip0.22cm
  \keybox{\textbf{建议路线：}短期保守声称 35T evidence-chain prototype；中期把 pointer semantic reconstruction 与 evasion-aware baseline 做成论文核心贡献。}
\end{frame}

% P23
\begin{frame}[plain]
  \begin{tikzpicture}[remember picture,overlay]
    \fill[accentcolor!8] (current page.north west) rectangle (current page.south east);
    \fill[accentcolor] (current page.north west) rectangle ([yshift=-0.16cm]current page.north east);
    \fill[accentcolor] ([yshift=0.16cm]current page.south west) rectangle (current page.south east);
    \node[text=accentcolor,font=\sffamily\bfseries\fontsize{26pt}{34pt}\selectfont]
      at ([yshift=0.35cm]current page.center) {请老师批评指正};
    \node[text=textgray,font=\sffamily\fontsize{10pt}{14pt}\selectfont,align=center]
      at ([yshift=-1.0cm]current page.center) {RV-MalTrace 论文进展汇报 \;|\; 2026年5月29日};
  \end{tikzpicture}
\end{frame}

\end{document}
